\section{Experiment Details}
\label{appx:experiment_details}

\subsection{Dataset Details}
\label{appx:dataset_details}

\begin{table}[thbp]
  \caption{Dataset descriptions. The dataset size is organized in (Train, Validation, Test).}\label{tab:dataset}
  \vskip 0.05in
  \centering
  \begin{threeparttable}
  \begin{small}
  \renewcommand{\multirowsetup}{\centering}
  \setlength{\tabcolsep}{3.8pt}
  \begin{tabular}{c|l|c|c|c|c}
    \toprule
    Tasks & Dataset & Dim & Series Length & Dataset Size & \scalebox{0.8}{Information (Frequency)} \\
    \toprule
     & ETTm1, ETTm2 & 7 & \scalebox{0.8}{\{96, 192, 336, 720\}} & (34465, 11521, 11521) & \scalebox{0.8}{Electricity (15 mins)}\\
    \cmidrule{2-6}
     & ETTh1, ETTh2 & 7 & \scalebox{0.8}{\{96, 192, 336, 720\}} & (8545, 2881, 2881) & \scalebox{0.8}{Electricity (Hourly)} \\
    \cmidrule{2-6}
    Long-term & Weather & 21 & \scalebox{0.8}{\{96, 192, 336, 720\}} & (36792, 5271, 10540) & \scalebox{0.8}{Weather (10 mins)} \\
    \cmidrule{2-6}
    Forecasting & Solar-Energy & 137 & \scalebox{0.8}{\{96, 192, 336, 720\}} & (36345, 5065, 10321) & \scalebox{0.8}{Solar (10 mins)} \\
    \cmidrule{2-6}
     & Exchange & 8 & \scalebox{0.8}{\{96, 192, 336, 720\}} & (5120, 665, 1422) & \scalebox{0.8}{Exchange rate (Daily)}\\
    \midrule
     & ETTm1, ETTm2 & 7 & 96 & (34465, 11521, 11521) & \scalebox{0.8}{Electricity (15 mins)} \\
    \cmidrule{2-6}
     & ETTh1, ETTh2 & 7 & 96 & (8545, 2881, 2881) & \scalebox{0.8}{Electricity (Hourly)}\\
    \cmidrule{2-6}
    Imputation & Electricity & 321 & 96 & (18317, 2633, 5261) & \scalebox{0.8}{Electricity (Hourly)}\\
    \cmidrule{2-6}
     & Weather & 21 & 96 & (36792, 5271, 10540) & \scalebox{0.8}{Weather (10 mins)} \\
    \cmidrule{2-6}
    & Exchange & 8 & 96 & (5120, 665, 1422) & \scalebox{0.8}{Exchange rate (Daily)}\\
    \midrule
     & \scalebox{0.8}{EthanolConcentration} & 3 & 1751 & (261, 0, 263) & \scalebox{0.8}{Alcohol Industry}\\
    \cmidrule{2-6}
    & \scalebox{0.8}{FaceDetection} & 144 & 62 & (5890, 0, 3524) & \scalebox{0.8}{Face (250Hz)}\\
    \cmidrule{2-6}
    & \scalebox{0.8}{Handwriting} & 3 & 152 & (150, 0, 850) & \scalebox{0.8}{Handwriting}\\
    \cmidrule{2-6}
    Classification & \scalebox{0.8}{Heartbeat} & 61 & 405 & (204, 0, 205)& \scalebox{0.8}{Heart Beat}\\
    \cmidrule{2-6}
    (UEA) & \scalebox{0.8}{JapaneseVowels} & 12 & 29 & (270, 0, 370) & \scalebox{0.8}{Voice}\\
    \cmidrule{2-6}
    & \scalebox{0.8}{SelfRegulationSCP1} & 6 & 896 & (268, 0, 293) & \scalebox{0.8}{Health (256Hz)}\\
    \cmidrule{2-6}
    & \scalebox{0.8}{SelfRegulationSCP2} & 7 & 1152 & (200, 0, 180) & \scalebox{0.8}{Health (256Hz)}\\
    \cmidrule{2-6}
    & \scalebox{0.8}{UWaveGestureLibrary} & 3 & 315 & (120, 0, 320) & \scalebox{0.8}{Gesture}\\
    \midrule
     & SMD & 38 & 100 & (566724, 141681, 708420) & \scalebox{0.8}{Server Machine} \\
    \cmidrule{2-6}
    Anomaly & MSL & 55 & 100 & (44653, 11664, 73729) & \scalebox{0.8}{Spacecraft} \\
    \cmidrule{2-6}
    Detection & SMAP & 25 & 100 & (108146, 27037, 427617) & \scalebox{0.8}{Spacecraft} \\
    \cmidrule{2-6}
     & SWaT & 51 & 100 & (396000, 99000, 449919) & \scalebox{0.8}{Infrastructure} \\
    \cmidrule{2-6}
    & PSM & 25 & 100 & (105984, 26497, 87841)& \scalebox{0.8}{Server Machine} \\
    \bottomrule
    \end{tabular}
    \end{small}
  \end{threeparttable}
\end{table}
As shown in Table \ref{tab:dataset}, the benchmark datasets span diverse domains and temporal granularities:
\begin{itemize}[leftmargin=*, itemsep=2pt]
    \item \textbf{Long-term Forecasting}: 7 datasets covering electricity transformer temperature (ETT series with 4 subsets), electricity consumption (Electricity), transportation flow (Traffic), meteorological indicators (Weather), solar power generation (Solar-Energy), and currency exchange rates (Exchange).
    \item \textbf{Imputation}: The same forecasting datasets with artificially induced random missing patterns at 12.5\%, 25\%, 37.5\%, and 50\% mask ratios, following the protocol established in \citep{wu2023timesnet}.
    \item \textbf{Classification}: 8 representative datasets from the UEA multivariate time series archive \citep{bagnall2018uea}, covering chemical sensing (EthanolConcentration), video analysis (FaceDetection), motion capture (Handwriting, UWaveGestureLibrary), medical diagnostics (Heartbeat), audio recognition (JapaneseVowels), and brain-computer interface (SelfRegulationSCP1, SelfRegulationSCP2).
    \item \textbf{Anomaly Detection}: 5 real-world datasets from server machines (SMD, PSM), spacecraft telemetry systems (MSL, SMAP), and water treatment infrastructure (SWaT).
\end{itemize}

\subsection{Baseline Methods}
\label{appx:baselines}
We reproduce all baselines using the official codebases and/or settings reported in the original papers.
To ensure a fair comparison, we keep the \emph{input embedding} and \emph{final projection} interfaces consistent across neural models when applicable, and focus on comparing their core sequence modeling or task-specific mechanisms.
Note that some baselines are only evaluated on specific tasks. %(see the task paragraphs below).

\paragraph{(1) Time-series specialized models.}
\begin{itemize}[leftmargin=*, itemsep=2pt]
\item \textbf{Transformer-based time-series models}:
\begin{itemize}[leftmargin=*, itemsep=1pt]
\item \textit{iTransformer} \citep{liu2023itransformer}: performs self-attention across variates (inverted tokenization) to better model multivariate dependencies.
\item \textit{PatchTST} \citep{nie2023patchtst}: uses patching and channel-independence to reduce sequence length and stabilize long-horizon learning.
\item \textit{Crossformer} \citep{zhang2023crossformer}: models cross-dimension dependencies via a two-stage attention design (intra-/inter-dimension).
\item \textit{Non-stationary Transformer} \citep{liu2022non} (reported as \textit{Stationary}): mitigates distribution shift with series stationarization and de-stationary attention.
\item \textit{Autoformer} \citep{wu2021autoformer}: replaces point-wise attention with auto-correlation and progressive series decomposition for long-term forecasting.
\item \textit{FEDformer} \citep{zhou2022fedformer}: conducts frequency-enhanced decomposition and attention in the spectral domain for efficient long-sequence modeling.
\item \textit{ETSformer} \citep{woo2022etsformer}: injects exponential smoothing (ETS-style) components to explicitly capture level/trend/seasonality.
\item \textit{Informer} \citep{zhou2021informer}: adopts ProbSparse attention to reduce attention cost for long sequence forecasting.
\item \textit{Reformer} \citep{kitaev2020reformer}: uses LSH attention to approximate softmax attention with lower memory/computation.
\item \textit{Pyraformer} \citep{liu2022pyraformer}: builds pyramidal attention across multiple temporal resolutions to model long-range structure.
\end{itemize}

\item \textbf{CNN/MLP-style time-series models}:
\begin{itemize}[leftmargin=*, itemsep=1pt]
\item \textit{TimesNet} \citep{wu2023timesnet}: reshapes 1D series into 2D representations guided by multi-periodicity and applies 2D convolution blocks.
\item \textit{DLinear} \citep{zeng2023dlinear}: decomposes series into trend and residual parts and forecasts with lightweight linear layers.
\item \textit{LightTS} \citep{campos2023lightts}: leverages efficient MLP mixing with interval/continuous sampling to handle long contexts.
\item \textit{TiDE} \citep{das2023long}: encodes the look-back with covariates via dense residual blocks and decodes future values with a per-time-step temporal decoder.
\item \textit{SCINet} \citep{liu2022scinet}: recursively splits and interacts subsequences to capture hierarchical temporal dependencies.
\item \textit{TCN} \citep{he2019temporal}: uses dilated causal convolutions to model long receptive fields without recurrence.
\end{itemize}

\item \textbf{RNN/State-space style models}:
\begin{itemize}[leftmargin=*, itemsep=1pt]
\item \textit{LSTM} \citep{hochreiter1997long}: recurrent gating mechanism for learning long/short-term temporal dependencies.
\item \textit{LSTNet} \citep{lai2018modeling}: combines CNN feature extraction with RNN temporal modeling and skip connections for periodic patterns.
\item \textit{LSSL} \citep{gu2022efficiently}: reparameterizes HiPPO-initialized state space models as diagonal-plus-low-rank, enabling near-linear-time long-sequence modeling via efficient Cauchy-kernel convolution..
\end{itemize}
\end{itemize}

\paragraph{(2) Pretrained LM-based baselines.}
\begin{itemize}[leftmargin=*, itemsep=2pt]
\item \textit{MAFS} \citep{huang2025many}: employs a multi-agent framework where each agent specializes in a forecasting sub-task, with inter-agent communication and a voting aggregator for final prediction. (Note: we include it here as it shares the pre-train then fine-tune paradigm, though it pre-trains iTransformer on sub-tasks rather than leveraging large language models.)
\item \textit{Time-VLM} \citep{zhong2025time}: integrates temporal, visual, and textual modalities through pretrained Vision-Language Models (e.g., ViLT, CLIP). Time series are transformed into images via frequency/periodicity encoding and multi-scale convolution, while contextual text prompts are generated from statistical features. A frozen VLM encodes both modalities, and cross-modal attention fuses them with temporal memory features for forecasting.
\item \textit{UniTime} \citep{liu2024unitime}: employs pretrained language models with domain instructions. It adopts channel-independence and patch-based tokenization, uses natural language instructions to provide domain identification and alleviate domain confusion, and applies masking to balance convergence speeds across domains. The Language-TS Transformer aligns time series from various input spaces to the common latent space of language models for cross-domain generalization.
\end{itemize}

\paragraph{(3) Classical baselines and others.}
\begin{itemize}[leftmargin=*, itemsep=1pt]
\item\textit{Transformer}\citep{vaswani2017attention}: the vanilla encoder--decoder attention architecture as a general-purpose sequence baseline.
\item \textit{DTW} \citep{Berndt1994UsingDT}: uses Dynamic Time Warping distance (typically with 1-NN) as a strong classical baseline for time-series classification.
\item \textit{XGBoost} \citep{Chen2016XGBoostAS}: gradient-boosted trees trained on engineered temporal features/statistics.
\end{itemize}
For anomaly detection, we additionally compare against \textit{Anomaly Transformer} \citep{xu2021anomaly}, a method specifically designed for this task that identifies anomalies via association discrepancy between series-wise and prior associations. Beyond task-specific models, we also consider classical methods that remain highly competitive on experimental results. For instance, in classification, \textit{DTW}, \textit{XGBoost}, and the vanilla \textit{Transformer} have all demonstrated strong performance and are thus included as additional baselines.

\subsection{Evaluation Metrics}

For the metrics, we adopt different evaluation criteria tailored to each task:

\paragraph{Long-term Forecasting and Imputation.} We use mean square error (MSE) and mean absolute error (MAE):
\begin{align*}
    \text{MSE} &= \frac{1}{H \cdot C} \sum_{i=1}^H \sum_{j=1}^C (\mathbf{X}_{i,j} - \widehat{\mathbf{X}}_{i,j})^2, &
    \text{MAE} &= \frac{1}{H \cdot C} \sum_{i=1}^H \sum_{j=1}^C |\mathbf{X}_{i,j} - \widehat{\mathbf{X}}_{i,j}|,
\end{align*}
where $\mathbf{X},\widehat{\mathbf{X}}\in\mathbb{R}^{H\times C}$ denote the ground truth and predicted values with $H$ time steps and $C$ feature dimensions. Lower values indicate better performance.


\paragraph{Classification.} We adopt accuracy (ACC) as the primary evaluation metric:
\begin{align*}
    \text{ACC} = \frac{\text{Number of Correct Predictions}}{\text{Total Number of Samples}} \times 100\%.
\end{align*}
Higher accuracy indicates better classification performance.


\paragraph{Anomaly Detection.} We adopt precision (P), recall (R), and F1-score (F1) following the point-adjust protocol \citep{xu2021anomaly}:
\begin{align*}
    \text{Precision} &= \frac{\text{TP}}{\text{TP} + \text{FP}}, &
    \text{Recall} &= \frac{\text{TP}}{\text{TP} + \text{FN}}, &
    \text{F1} &= 2 \cdot \frac{\text{Precision} \cdot \text{Recall}}{\text{Precision} + \text{Recall}},
\end{align*}
where TP, FP, and FN denote true positives, false positives, and false negatives, respectively. The F1-score provides a balanced measure between precision and recall, with higher values indicating better detection performance.

\subsection{Model and Agent Configuration}
\label{appx:model_agent_config}

This section provides a comprehensive overview of the model hyperparameters, agent-specific configurations, and shared memory mechanism in \model. The framework adopts a modular design where each agent maintains specialized parameter sets while coordinating through a unified memory interface.

\paragraph{Training Configuration.}
Table~\ref{tab:model_config} summarizes the task-specific training configurations. All experiments employ the AdamW optimizer with default momentum parameters $(\beta_1, \beta_2) = (0.9, 0.999)$ and weight decay of $10^{-2}$. We apply early stopping with a patience of 3 epochs to prevent overfitting and use cosine annealing for learning rate scheduling. Mixed-precision training (FP16) is enabled by default to accelerate computation while maintaining numerical stability.

\begin{table}[thbp]
  \vspace{-5pt}
  \caption{Task-specific training configurations for \model. The embedding dimension $d_{\text{model}}$ and network depth are kept consistent across tasks to ensure fair comparison, while learning rates and batch sizes are tuned per task for optimal convergence.}\label{tab:model_config}
  \vskip 0.05in
  \centering
  \begin{threeparttable}
  \begin{small}
  \renewcommand{\multirowsetup}{\centering}
  \setlength{\tabcolsep}{3.8pt}
  \begin{tabular}{l|cccc|cccc}
    \toprule
    \multirow{2}{*}{\textbf{Task}} & \multicolumn{4}{c|}{\textbf{Architecture}} & \multicolumn{4}{c}{\textbf{Training}} \\
    \cmidrule(lr){2-5}\cmidrule(lr){6-9}
    & $d_{\text{model}}$ & Layers & Patch & Stride & LR & Loss & Batch & Epochs\\
    \midrule
    Long-term Forecast & 64 & 2 & 16 & 8 & $10^{-4}$ & MSE & 32 & 10 \\
    Imputation & 64 & 2 & 16 & 8 & $10^{-3}$ & MSE & 16 & 10 \\
    Classification & 64 & 2 & 16 & 8 & $10^{-3}$ & CE & 16 & 30 \\
    Anomaly Detection & 64 & 3 & 16 & 8 & $10^{-4}$ & MSE & 128 & 10 \\
    \bottomrule
  \end{tabular}
  \begin{tablenotes}
    \footnotesize
    \item[] LR: initial learning rate with cosine annealing decay. CE: Cross-Entropy loss.
  \end{tablenotes}
  \end{small}
  \end{threeparttable}
  \vspace{-5pt}
\end{table}

\paragraph{Ablation Flags.}
\model supports fine-grained ablation studies through configurable flags that enable/disable individual components:
\begin{itemize}[leftmargin=*, itemsep=1pt]
    \item \texttt{enable\_visual\_reasoner}: Toggle VLM-based anchor extraction (default: True)
    \item \texttt{enable\_numeric\_reasoner}: Toggle Neural ODE sequence reconstruction (default: True)
    \item \texttt{enable\_shared\_memory}: Toggle cross-agent memory sharing (default: True)
    \item \texttt{enable\_gated\_attention}: Toggle gated vs. simple MLP communication (default: True)
    \item \texttt{enable\_tools}: Toggle tool-augmented execution vs. simple MLP heads (default: True)
    \item \texttt{completion\_strategy}: ODE solver alternative: \{`ode', `linear', `quadratic', `repeat'\}
\end{itemize}

\begin{table}[htbp]
  \centering
  \caption{MAS4TS Model Architecture Parameters}
  \label{tab:model_params}
  \begin{small}
    \scalebox{0.95}{
      \begin{tabular}{l|c|p{5.8cm}}
        \toprule
        \textbf{Parameter} & \textbf{Value} & \textbf{Description} \\
        \midrule
        \multicolumn{3}{l}{\textit{Global Configuration}} \\
        \midrule
        d\_model & 64 & Embedding dimension for all agents \\
        d\_memory & 64 & Shared memory dimension \\
        patch\_len & 16 & Patch length for tokenization \\
        stride & 8 & Stride between patches \\
        \midrule
        \multicolumn{3}{l}{\textit{Data Analyzer Agent}} \\
        \midrule
        top\_k\_features & 10 & Top-k features for covariance selection \\
        feature\_method & covariance & Feature selection criterion \\
        \midrule
        \multicolumn{3}{l}{\textit{Visual Reasoner Agent}} \\
        \midrule
        vlm\_model & Qwen3-VL-235B-Instruct & Vision-Language Model backbone \\
        max\_anchors & 20 & Maximum anchor points per sample \\
        confidence\_threshold & 0.7 & Minimum confidence for anchor acceptance \\
        \midrule
        \multicolumn{3}{l}{\textit{Numeric Reasoner Agent}} \\
        \midrule
        hidden\_dim & 128 & Hidden dimension for ODE dynamics \\
        ode\_solver & RK4 & ODE integration method \\
        ode\_step & 0.05 & Integration step size \\
        \midrule
        \multicolumn{3}{l}{\textit{Task Executor Agent}} \\
        \midrule
        e\_layers & 2 & Encoder layers in task tools \\
        n\_heads & 4 & Attention heads \\
        dropout & 0.1 & Dropout rate \\
        \bottomrule
      \end{tabular}
    }
  \end{small}
\end{table}


\subsection{Implementation Details}
We provide the dataset descriptions and experiment configurations in Table~\ref{tab:dataset} and~\ref{tab:model_config}. All experiments are conducted on a single NVIDIA GeForce RTX 4090 GPU (48GB version) with a fixed random seed for reproducibility. The MAS4TS framework is implemented in PyTorch 2.9.1+cu128 with CUDA 12.8. For the Vision-Language Model component, we utilize the \textit{Qwen2-VL-235B-Instruct} model accessed through EAS API services of the Alibaba Cloud Platform. 